\SourceLeaf{249}%
\noindent considered in abstracto, involves no contradiction and might well be accessible to higher beings. But to return from this height to earth, it is a pity that the equality of volumes of solid figures that are merely symmetric and lack congruence can be demonstrated only by the method of exhaustion, and cannot be demonstrated as elementarily as you, to my knowledge, were the first to show in the case of the area of the spherical triangle. Your proof occurs in a somewhat different form in \textsc{Lobatschewsky}, though it is unclear whether he knew your paper. Every spherical triangle can be divided into three isosceles triangles. Applied to two symmetric triangles, this division gives equal triangles in both, differing only in their arrangement. (It goes without saying that for completeness one would have to distinguish three cases, of which the second and third can also be stated by saying that one of the component triangles in each case may be negative or 0.)
\mbox{.\ .\ .\ .\ .\ .\ .\ .}

\par\bigskip

\Needspace{7\baselineskip}
\begin{center}
[\textsc{Gerling} \emph{to} \textsc{Gauss}.\quad \emph{Marburg, 15 April 1844.}]
\end{center}

\{Although I cannot properly answer your letter of the 8th of this month, for which I offer my warm thanks in the meantime, because I have yet to find time to study it and absorb it as fully as its rich contents and my interest in the subject demand, I cannot deny myself the pleasure of writing you a few lines about a particular remark in it.

You say: “It is a pity that the equality of volumes of solid figures that are merely symmetric and lack congruence can be demonstrated only by the method of exhaustion, and cannot be demonstrated as elementarily”
etc.

Allow me to communicate such a demonstration here and to ask you to write to me on some occasion with your judgment of it.

\textls{Theorem}: Any two symmetric polyhedra have equal volumes.

\textls{Proof}: I. By drawing suitable planes through them, they can be divided\space%
